\subsection{显式算法与复杂度界（审稿级）}\label{subsec:zeckendorf-algorithms-complexity}

本小节把“可计算、可终止”的表述提升为\textbf{显式算法}与\textbf{复杂度界}。为保持依赖最小，算法采用最短闭环：\textbf{值层运算}（整数加法/乘法）+\textbf{Zeckendorf 规范化}（贪心）。

\subsubsection{Zeckendorf 规范化：从整数到合法数字}

给定 $N\in\NN$，目标是计算其唯一 Zeckendorf 数字 $Z(N)=(c_1,c_2,\dots)$（仅有限个 $1$ 且无相邻 $1$）使得 $N=\sum_{k\ge 1}c_kF_{k+1}$。

\begin{itemize}
  \item \textbf{输入：}整数 $N$。
  \item \textbf{输出：}Zeckendorf 合法数字序列 $Z(N)$。
\end{itemize}

\textbf{算法（贪心 + 跳邻位）：}
\begin{enumerate}
  \item 预生成 Fibonacci 权重 $F_{2},F_{3},\dots,F_{M+1}$，其中 $M$ 足够大使得 $F_{M+1}\le N<F_{M+2}$；
  \item 从 $k=M$ 递减到 $1$：若 $F_{k+1}\le N$，则置 $c_k:=1$、令 $N\leftarrow N-F_{k+1}$，并跳过相邻位（将 $k\leftarrow k-1$ 再继续递减）；否则置 $c_k:=0$；
  \item 输出 $c=(c_1,c_2,\dots,c_M)$（其后全为 $0$）。
\end{enumerate}

该算法的正确性与唯一性由 Zeckendorf 定理给出（\cite{Zeckendorf1972}；本稿中见附录 C）。

\subsubsection{稳定加法与稳定乘法：两步实现}

对输入 $c,d\in X_\infty$（仅有限个 $1$），令其有效分辨率为
\[
m:=\max\{k:\ c_k=1\text{ 或 }d_k=1\}.
\]
则
\begin{itemize}
  \item \textbf{稳定加法} $c\oplus d$：计算 $N=\mathrm{Val}(c)+\mathrm{Val}(d)$，再输出 $Z(N)$；
  \item \textbf{稳定乘法} $c\otimes d$：计算 $N=\mathrm{Val}(c)\mathrm{Val}(d)$，再输出 $Z(N)$。
\end{itemize}
这正是附录 E 的“值层运算 + 规范形”的最短闭环实现。

\subsubsection{复杂度（按分辨率 \texorpdfstring{$m$}{m} 与按比特两种口径）}

为了避免模型争议，本文在附录 E$'$ 中同时给出两种等价口径的复杂度界：分辨率口径（把 $m$ 视为输入长度）与比特口径（把整数按二进制长度计成本）。结论可概括如下（详证见附录 E$'$ 中的命题 \ref{prop:fold-complexity} 与 \ref{prop:stable-arith-complexity}）：

\begin{itemize}
  \item \textbf{折叠/规范化：}对长度 $m$ 的输入，Zeckendorf 规范化在分辨率口径为 $O(m)$ 步；在比特口径为 $O(m^2)$ 比特运算。
  \item \textbf{稳定加法：}分辨率口径 $O(m)$；比特口径 $O(m^2)$。
  \item \textbf{稳定乘法：}分辨率口径为“值层乘法成本 + $O(m)$”；比特口径为 $O(M(m)+m^2)$，其中 $M(m)$ 为 $\\Theta(m)$ 比特整数乘法成本（学校乘法为 $O(m^2)$）。
\end{itemize}

\subsubsection{纯局部归一化（可选扩展）}\label{subsec:local-normalize}
本方向用 Zeckendorf 数字上的局部 carry/normalize 直接实现 $\Fold_\infty$，避免显式值层运算。以逐位叠加得到的串 $s\in\{0,1,2\}^m$ 为输入，给出一个三趟局部规则体系（每步保持 $\mathrm{Val}$；规则可由 Fibonacci 递推验证，见 \cite{AhlbachUsatineFrougnyPippenger2013}）：
\begin{itemize}
  \item \textbf{第 1 趟（宽度 4，左到右）}：对 $x\in\{0,1,2\}$，
  \[
  020x\Rightarrow 100x^{+},\quad
  030x\Rightarrow 110x^{+},\quad
  021x\Rightarrow 110x,\quad
  012x\Rightarrow 101x,
  \]
  其中 $x^{+}:=x+1$。末端清理可取
  \[
  030\Rightarrow 111,\quad 020\Rightarrow 101,\quad 0120\Rightarrow 1010,\quad
  03\Rightarrow 11,\quad 02\Rightarrow 10,\quad 012\Rightarrow 101.
  \]
  \item \textbf{第 2--3 趟（宽度 3）}：用规则
  \[
  011\Rightarrow 100
  \]
  先右到左、再左到右各扫一遍，消去相邻 $11$。
\end{itemize}

\begin{theorem}[三趟局部归一化与 $\Fold_\infty$ 一致]\label{thm:local-normalize-fold}
令 $T:=T_3\circ T_2\circ T_1$ 为上述三趟归一化。则对任意 $s\in\{0,1,2\}^m$，
\[
T(s)=\Fold_\infty(s)\in X_m.
\]
\end{theorem}

\textbf{证明：}
每条规则保持 $\mathrm{Val}$（等价地：在群胚生成元轨道内移动），且第 2--3 趟保证输出无相邻 $11$，从而落在 $X_m$。由定理 \ref{thm:zeckendorf-transversal} 的横截面唯一性交点，终止形必为 $\Fold_\infty(s)$。证毕。

\begin{corollary}[有限状态实现与线性时间]\label{cor:local-normalize-fst}
上述三趟均可由有限状态转导器以常数延迟实现，故分辨率口径下归一化为 $O(m)$。例如第 1 趟可取状态集 $Q_1=\{0,1,2,3\}^3$（至多 $4^3$ 个状态），第 2--3 趟可取 $Q_2=\{0,1\}^2$（至多 $2^2$ 个状态）。\cite{AhlbachUsatineFrougnyPippenger2013,Frougny1999OnlineAddition}
\end{corollary}

\subsubsection{单趟在线延迟 \texorpdfstring{$3$}{3}：统一转导器}\label{subsec:online-delay}
三趟局部规则可视为转导器串联的产品自动机。更强的结果是存在\textbf{单趟在线}转导器（延迟 3）直接实现同一归一化；其同步部分可用“余量守恒”给出（\cite{Frougny1999OnlineAddition}）。

令 $\varphi=(1+\sqrt5)/2$，定义余量
\[
R(s):=s_1\varphi^{-1}+s_2\varphi^{-2}+s_3\varphi^{-3}\in[0,1),
\]
并令输入 $d\in\{0,1,2\}$、输出 $e\in\{0,1\}$。同步边 $(s\to t)$ 满足守恒式
\[
R(s)+d\varphi^{-4}=e\varphi^{-1}+t_1\varphi^{-2}+t_2\varphi^{-3}+t_3\varphi^{-4},
\]
等价地
\[
R(t)=\varphi R(s)+d\varphi^{-3}-e,\qquad R(t)\in[0,1).
\]
\begin{remark}[延迟-$3$ 归一化的扩张--取整表示]
由 $R(t)\in[0,1)$ 可写成
\[
e=\Bigl\lfloor \varphi R(s)+d\varphi^{-3}\Bigr\rfloor,\qquad
R(t)=\Bigl\{\varphi R(s)+d\varphi^{-3}\Bigr\},
\]
其中 $\{\cdot\}$ 为小数部分。由于同步核只允许 $e\in\{0,1\}$（附录 B 的转移表中无 $e=2$ 边），上述在线核等价于在 $[0,1)$ 上由扩张映射 $r\mapsto \varphi r+d\varphi^{-3}$ 生成、并以取整 cocycle 回投影的动力系统。该表示把“在线延迟”直接编码为乘法扩张与分支选择的串行结构。
\end{remark}
该同步部分的状态集可取为
\[
Q=\{000,001,002,010,100,101,0\overline{1}2,1\overline{1}2,01\overline{1},11\overline{1}\},
\]
并且出边（标记 $d/e$）可写为：
\[
\begin{aligned}
000&\xrightarrow{d/0}00d,\qquad 100\xrightarrow{d/1}00d\quad(d=0,1,2),\\
001&\xrightarrow{0/0}010,\ \xrightarrow{1/0}100,\ \xrightarrow{2/0}101,\\
002&\xrightarrow{0/0}11\overline{1},\ \xrightarrow{1/1}000,\ \xrightarrow{2/1}001,\\
010&\xrightarrow{0/0}100,\ \xrightarrow{1/0}101,\ \xrightarrow{2/1}0\overline{1}2,\\
101&\xrightarrow{0/1}010,\ \xrightarrow{1/1}100,\ \xrightarrow{2/1}101,\\
0\overline{1}2&\xrightarrow{0/0}01\overline{1},\ \xrightarrow{1/0}010,\ \xrightarrow{2/0}100,\\
1\overline{1}2&\xrightarrow{0/1}01\overline{1},\ \xrightarrow{1/1}010,\ \xrightarrow{2/1}100,\\
01\overline{1}&\xrightarrow{0/0}001,\ \xrightarrow{1/0}002,\ \xrightarrow{2/0}1\overline{1}2,\\
11\overline{1}&\xrightarrow{0/1}001,\ \xrightarrow{1/1}002,\ \xrightarrow{2/1}1\overline{1}2.
\end{aligned}
\]
末端函数可取
\[
\omega(s_1s_2s_3):=Z(s_1F_4+s_2F_3+s_3F_2),
\]
以吞吐延迟窗口的残余并落回 Zeckendorf 规范形（\cite{Frougny1999OnlineAddition}）。

\begin{theorem}[单趟在线延迟 3 的 $\Fold_\infty$ 实现]\label{thm:online-delay-fold}
以上在线转导器（含末端函数）对任意 $s\in\{0,1,2\}^m$ 输出 $\Fold_\infty(s)$；其同步图即附录 B 的 10 状态同步核。
\end{theorem}

\textbf{证明：}
守恒式保证值不变量，末端函数输出落在 $X_m$；由定理 \ref{thm:zeckendorf-transversal} 的横截面唯一性得 $\Fold_\infty(s)$。同步图与附录 B 的状态/转移一致。证毕。

\begin{remark}[操作数不一致与隐含属性]
原子生成元 $E_k,E_k^{-1}$ 具有尺度参数 $k$ 与守卫条件，体现为“隐含类型”。令 $\nu(a):=\sum_k a_k$ 表示块数，则
\[
\nu(E_k(a))=\nu(a)+1,\qquad \nu(E_k^{-1}(a))=\nu(a)-1,
\]
即局部裂解/合并天然呈现 $1\leftrightarrow 2$ 的操作数不一致性。
\end{remark}

\begin{remark}[primitive 轨道与“素数影子”]\label{rem:primitive-shadow}
上述在线转导器的同步图是有限有向图，其 primitive 周期轨道给出“不可约循环模式”。设其邻接矩阵为 $B$，则周期点计数 $a_n=\mathrm{Tr}(B^n)$，primitive 轨道数 $p_n$ 由
\[
a_n=\sum_{d\mid n} d\,p_d,\qquad
p_n=\frac{1}{n}\sum_{d\mid n}\mu(d)\,a_{n/d}
\]
给出，并与第 9 节的 $\zeta$--欧拉乘积结构一致（参见 \cite{LindMarcus1995} 与附录 B 中的 $\zeta_B$、显式分解与数值谱）。
\end{remark}

\begin{remark}[素数影子与算术素数的层级差异]\label{rem:prime-shadow-vs-primes}
有限核的 primitive 轨道谱是动力系统层的 prime-like 对象，严格由该核的 $\zeta$ 与谱决定；它不等同于整数素数集合本身。任何固定分辨率/有限状态的读出只能携带有限信息，因此无法在语义上与“整数素数”等价。若要把“素数”作为本体结构引入，需在更高层把素数作为算子族的原子标签：例如在 Hecke 对应框架中以 $T_p$ 作为生成算子，并在本征基上读出 $a_p$；此时素数寄存器是算子寄存器而非因子分解寄存器。本文的同步核谱与 $\zeta$ 计算提供的是动力学模板层，可作为这类升级接口的底座。
\end{remark}

\subsubsection{时间/空间对偶的交换图：trace--iterate 与 Euler product 的严格一致}\label{subsec:time-space-commuting-diagram}
本节把“一个像时间（迭代/迹），一个像空间（primitive/欧拉乘积）”的直觉提升为严格的交换图。核心点是：在有限矩阵（或迹类算子）层面，$\det(I-zT)$、$\mathrm{Tr}(T^n)$ 与 primitive 抽取（Möbius/Witt）都为共轭不变量，因此可以函子化地组织为一张可审计的交换图。

\paragraph{三个范畴（只保留共轭不变量）}
取如下三类对象：
\begin{itemize}
  \item \textbf{端态范畴 $\mathbf{End}^{\simeq}$}：对象为对 $(V,T)$（有限维线性空间与算子 $T:V\to V$），态射为相似变换（共轭）；
  \item \textbf{序列范畴 $\mathbf{Seq}$}：对象为序列 $a=(a_n)_{n\ge 1}$（取值于 $\CC$），取离散范畴口径；
  \item \textbf{Euler 范畴 $\mathbf{Euler}$}：对象为形式幂级数 $F(z)\in 1+z\CC[[z]]$（取离散范畴口径）。
\end{itemize}

\paragraph{四个函子（时间侧、原子化、Euler product 与行列式 $\zeta$）}
定义：
\[
\mathcal{T}:\mathbf{End}^{\simeq}\to\mathbf{Seq},\qquad
\mathcal{T}(T)=(a_n)_{n\ge 1},\ a_n:=\mathrm{Tr}(T^n),
\]
\[
\mathcal{M}:\mathbf{Seq}\to\mathbf{Seq},\qquad
(\mathcal{M}a)_n:=\frac{1}{n}\sum_{d\mid n}\mu(d)\,a_{n/d},
\]
\[
\mathcal{E}:\mathbf{Seq}\to\mathbf{Euler},\qquad
\mathcal{E}(p):=\prod_{n\ge 1}(1-z^n)^{-p_n},
\]
\[
\mathcal{Z}:\mathbf{End}^{\simeq}\to\mathbf{Euler},\qquad
\mathcal{Z}(T):=\frac{1}{\det(I-zT)}=\exp\!\left(\sum_{n\ge 1}\frac{\mathrm{Tr}(T^n)}{n}z^n\right).
\]

\begin{proposition}[交换图（严格交换）]\label{prop:time-space-commuting-diagram}
在上述记号下有恒等式
\[
\boxed{\ \mathcal{Z}=\mathcal{E}\circ\mathcal{M}\circ\mathcal{T}\ }.
\]
等价地，令 $a_n=\mathrm{Tr}(T^n)$，令 $p=\mathcal{M}(a)$，则
\[
\frac{1}{\det(I-zT)}=\prod_{n\ge 1}(1-z^n)^{-p_n}\in 1+z\CC[[z]].
\]
\end{proposition}

\begin{proof}
由定义 $\log\mathcal{E}(p)=-\sum_{n\ge 1}p_n\log(1-z^n)$，展开得
\[
\log\mathcal{E}(p)=\sum_{n\ge 1}p_n\sum_{k\ge 1}\frac{z^{nk}}{k}
=\sum_{m\ge 1}\left(\sum_{n\mid m}\frac{n\,p_n}{m}\right)z^m.
\]
另一方面，$p=\mathcal{M}(a)$ 等价于 $a_m=\sum_{n\mid m}n\,p_n$（周期闭环分解为 primitive 幂重复），故
\(
\log\mathcal{E}(p)=\sum_{m\ge 1}\frac{a_m}{m}z^m
\)。
取指数即得 $\mathcal{E}(p)=\exp(\sum_{m\ge 1}\frac{\mathrm{Tr}(T^m)}{m}z^m)=\mathcal{Z}(T)$。证毕。
\end{proof}

\begin{remark}[同步核的自对偶与二变量函数方程]
取输出位势 $\varphi(e)=e$，令 $B(u)=B_0+uB_1$、$\Delta(z,u)=\det(I-zB(u))$ 与 $\zeta(z,u)=\Delta(z,u)^{-1}$。附录 \ref{app:kernel-self-dual} 给出置换矩阵 $\Pi$ 使得 $\Pi^{-1}B(u)\Pi=uB(1/u)$，因此
\[
\Delta(z,u)=\Delta(uz,1/u),\qquad
\zeta(z,u)=\zeta(uz,1/u).
\]
这可视作端态侧对偶函子 $D_{\mathbf{End}}:B(u)\mapsto uB(1/u)$ 与 Euler 侧对偶函子 $D_{\mathbf{Euler}}:F(z,u)\mapsto F(uz,1/u)$ 之间的严格交换，从而把“时间长度 $z$”与“事件权重 $u$”在同一 $\zeta$-接口下对齐为可核验的对偶变换。
\end{remark}

\begin{remark}[完成化与 \texorpdfstring{$\ZZ/2$}{Z2} 作用]
把 $D_{\mathbf{End}}$ 与 $D_{\mathbf{Euler}}$ 视为对参数 $u$ 的对偶作用，则 $D^2=\mathrm{Id}$ 给出一条 $\ZZ/2$-对称。若 $F(z,u)$ 满足 $F(z,u)=F(uz,1/u)$，定义完成化
\[
\widetilde{F}(z,\theta):=F(e^{-\theta/2}z,e^\theta),\qquad u=e^\theta,
\]
则 $\widetilde{F}(z,-\theta)=\widetilde{F}(z,\theta)$ 为严格偶性。对应严格交换图
\[
\begin{CD}
\mathbf{End}^{\simeq} @>D_{\mathbf{End}}>> \mathbf{End}^{\simeq}\\
@V\mathcal{Z}VV @VV\mathcal{Z}V\\
\mathbf{Euler} @>D_{\mathbf{Euler}}>> \mathbf{Euler}
\end{CD}
\]
写作完成化算子
\[
\mathcal{C}(F)(z,\theta):=F(e^{-\theta/2}z,e^\theta),
\]
则 $\mathcal{C}\circ D_{\mathbf{Euler}}=\mathcal{C}$，并且
\[
\widetilde{\mathcal{Z}}:=\mathcal{C}\circ\mathcal{Z}
\]
满足 $\widetilde{\mathcal{Z}}=\widetilde{\mathcal{Z}}\circ D_{\mathbf{End}}$。因此对同步核的 $\Delta$ 与 $\zeta$，完成化后在 $\theta$ 上呈严格偶性（见附录 \ref{app:kernel-completion}）。
\end{remark}

\begin{remark}[为何 delay-$3$ 在线机是“prime-like 原子的产生机”]
定理 \ref{thm:online-delay-fold} 给出一个具体对象：delay-$3$ 在线归一化同步核的邻接矩阵 $B$（附录 B）。将 $T=B$ 代入命题 \ref{prop:time-space-commuting-diagram}，则“时间侧”的 $\mathrm{Tr}(B^n)$ 与“空间侧”的 primitive 轨道数 $p_n$（prime-like 原子）严格由同一 $\zeta_B(z)=1/\det(I-zB)$ 统一。
因此“不可并行依赖”（延迟窗口中的余量守恒/携带）并非噪声，而是把迭代动力学编码为一个有限核，使得其 primitive 闭环成为可审计的 prime-like 原子。
\end{remark}

\begin{remark}[乘法相关性的可检验落点：张量装配与 Witt/Möbius]
当把单流核按 $m$ 条同构同余流做全局块级张量装配时，有 $T_{\mathrm{glob}}=T_{\mathrm{flow}}^{\otimes m}$，从而
\[
a_n^{\mathrm{glob}}(u):=\mathrm{Tr}(T_{\mathrm{glob}}(u)^n)=\bigl(a_n^{\mathrm{flow}}(u)\bigr)^m.
\]
在附录 \ref{app:witt-weighted} 的 $u$-加权口径下，primitive 谱多项式满足
\[
p_n^{\mathrm{glob}}(u)=\frac{1}{n}\sum_{d\mid n}\mu(d)\,a_{n/d}^{\mathrm{glob}}(u^d)
=\frac{1}{n}\sum_{d\mid n}\mu(d)\,\bigl(a_{n/d}^{\mathrm{flow}}(u^d)\bigr)^m.
\]
特别地，若 $n=\ell$ 为素数，则
\[
p_\ell^{\mathrm{glob}}(u)=\frac{1}{\ell}\left(\bigl(a_\ell^{\mathrm{flow}}(u)\bigr)^m-\bigl(a_1^{\mathrm{flow}}(u^\ell)\bigr)^m\right).
\]
该等式把“时间侧的幂次/乘法”精确翻译为“空间侧 primitive 谱的非平凡 Witt 多项式律”（见推论 \ref{cor:par-add-global-lift-endpoints} 的同型口径），从而给出“不可并行活动”的一个可审计 prime-like 读出接口。
\end{remark}

